\documentclass[11pt,a4paper]{article}

\usepackage[T1]{fontenc}
\usepackage[utf8]{inputenc}
\usepackage{mathptmx}
\usepackage{amsmath,amssymb,amsthm}
\usepackage{geometry}
\usepackage{microtype}
\usepackage{hyperref}
\usepackage{booktabs}
\usepackage{enumitem}
\usepackage{setspace}
\usepackage{titlesec}
\usepackage{abstract}
\usepackage{natbib}
\usepackage{xcolor}

\geometry{top=1.1in, bottom=1.1in, left=1.15in, right=1.15in}

\titleformat{\section}{\normalfont\bfseries\large}{\thesection}{1em}{}
\titleformat{\subsection}{\normalfont\bfseries\normalsize}{\thesubsection}{1em}{}
\titleformat{\subsubsection}{\normalfont\bfseries\itshape\normalsize}{\thesubsubsection}{1em}{}

\theoremstyle{plain}
\newtheorem{theorem}{Theorem}[section]
\newtheorem{proposition}[theorem]{Proposition}
\newtheorem{corollary}[theorem]{Corollary}
\theoremstyle{definition}
\newtheorem{definition}[theorem]{Definition}
\theoremstyle{remark}
\newtheorem{remark}[theorem]{Remark}

\newcommand{\ARA}{A_{\mathrm{RA}}}
\newcommand{\tARA}{\tilde{A}_{\mathrm{RA}}}
\newcommand{\TU}{T_U}

\hypersetup{
  colorlinks=true, linkcolor=blue, citecolor=blue, urlcolor=blue,
  pdftitle={The Causal Firewall: Substrate-Independent Physical Constraints on the Distribution of Life},
  pdfauthor={Joshua F. Sandeman}
}

\begin{document}

% ── Title block ───────────────────────────────────────────────────────────────
\begin{center}
  {\LARGE\bfseries The Causal Firewall:\\ Substrate-Independent Physical Constraints\\
  on the Distribution of Life in the Universe\par}
  \vspace{14pt}
  {\large Joshua F.\ Sandeman$^{1*}$\par}
  \vspace{4pt}
  {\normalsize $^{1}$Independent Researcher, Salem, OR, USA.\par}
  {\normalsize $^{*}$sansuikyo@gmail.com\par}
  \vspace{8pt}
  {\small	extit{Preprint available at}~\href{https://doi.org/10.5281/zenodo.19319374}{doi.org/10.5281/zenodo.19319374}\par}
  \vspace{14pt}
\end{center}

% ── Abstract ─────────────────────────────────────────────────────────────────
\begin{abstract}
We derive a substrate-independent physical criterion for the distribution
of biological complexity in the universe. The criterion, called the
\emph{Causal Firewall}, follows from a single physical requirement: that
the environment maintain a universally agreed-upon causal ledger of
actualization events over the timescales required for complex molecular
assembly. Two conditions are necessary: (F1) the Quantum Darwinism
redundancy threshold, requiring that actualized states are encoded
in at least $R_{\min} \sim \mathcal{O}(10)$ independent environmental
fragments; and (F2) the non-relativistic decoherence condition, requiring
that the local Unruh temperature $T_U = \hbar a / (2\pi c k_B)$ satisfy
$T_U \ll T_{\mathrm{env}}$. Crucially, neither condition makes reference
to molecular class, solvent chemistry, temperature range, or biological
alphabet. We show that the Causal Firewall corresponds to the $\mu = 1$
Erd\H{o}s-R\'{e}nyi percolation threshold of the local causal graph ---
the same mathematical transition governing QCD confinement, galactic
rotation curves, and quantum fault-tolerance thresholds. We further derive
a quantitative constraint on origin-of-life scenarios: the coarse-grained
assembly depth $\tilde{A}_{\mathrm{RA}}$ of the first proto-metabolic
network is bounded both above (by the Causal Firewall cycle time constraint)
and below (by the Assembly Theory floor). The framework predicts that life
is structurally impossible in highly relativistic environments regardless
of energy availability, that non-carbon-based life is not prohibited but
is subject to the same causal coherence conditions, and that SETI search
programmes should be weighted toward environments satisfying the Causal
Firewall criteria rather than merely toward circumstellar habitable zones.
Two falsifiable predictions are stated.
\end{abstract}

\noindent\textbf{Keywords:} astrobiology; habitability; causal set theory;
Unruh effect; quantum decoherence; origin of life; SETI; Assembly Theory

\medskip
\noindent\textbf{Ethics statement.}
This article does not contain any studies with human participants or animals.

% ── 1. Introduction ──────────────────────────────────────────────────────────
\section{Introduction}
\label{sec:intro}

The circumstellar habitable zone (CHZ) --- the annular region around a star
where liquid water can exist on a planetary surface --- has been the dominant
framework for identifying environments conducive to life since its modern
formulation by \citet{Kasting1993} and its refinement by \citet{Kopparapu2013}.
Despite its practical utility as a search criterion for biosignatures, the
CHZ framework rests on an implicit assumption: that terrestrial biochemistry
is universal, or at least that liquid water is a necessary condition for
life everywhere. This assumption is a proxy, not a derivation.

Alternative biochemistry proposals --- life in liquid methane, ammonia solvents,
or non-aqueous environments \citep{Benner2004} --- challenge the water-centrism
of the CHZ without replacing it with a more fundamental criterion. They expand
the list of permitted solvents while leaving the underlying question unanswered:
what is the \emph{physical reason} that any environment either permits or
prohibits biological complexity, independently of its specific chemistry?

Here we derive such a criterion from first principles. The Relational Actualism
(RA) framework \citep{Sandeman2026RACI} grounds biological complexity in the
irreversibility of physical interactions: an organism's present state is not
merely correlated with its past, it is \emph{constituted} by the irreversible
sequence of on-shell actualization events that produced it
\citep{Sandeman2026RAQM}. This ontology makes the emergence of life a physical
question about causal coherence rather than a chemical question about molecular
structure. From it follows a set of necessary physical conditions for biological
complexity --- conditions we call the \textbf{Causal Firewall} --- that make
no reference to molecular class, solvent, or temperature.

We show that the Causal Firewall corresponds to a percolation threshold in the
local causal graph, that it is the same mathematical transition appearing at
four other scales in the RA framework, and that it yields three quantitative
astrophysical observables measurable for any candidate environment. We derive
a novel ``sandwich bound'' on the causal depth of the first self-replicating
systems and state two falsifiable predictions for the distribution of life.

\paragraph{Structure.}
Section~\ref{sec:firewall} presents the Causal Firewall conditions and their
substrate-independence. Section~\ref{sec:percolation} connects the Firewall
to the Erd\H{o}s-R\'{e}nyi transition and the universal $\mu = 1$ phenomenon.
Section~\ref{sec:ara} derives the coarse-grained assembly depth and the
origin-of-life sandwich bound. Section~\ref{sec:predictions} develops three
universal biosignature criteria and two falsifiable predictions.
Section~\ref{sec:comparison} situates the framework relative to existing
habitability theories. Section~\ref{sec:conclusion} concludes.

% ── 2. The Causal Firewall ────────────────────────────────────────────────────
\section{The Causal Firewall}
\label{sec:firewall}

\subsection{Causal Ledger Preservation}

In the RA framework, a biological organism is a physical system whose
complexity just \emph{is} its actualization history: the unbroken sequence
of irreversible on-shell interactions that assembled it, layer by layer,
from elementary constituents \citep{Sandeman2026RAHC}. Each irreversible
interaction writes a permanent directed edge into the growing causal Directed
Acyclic Graph (DAG) $\mathcal{G} = (\mathcal{V}, \mathcal{E})$. The assembly
of any complex molecule requires a specific sequential ordering of these
interactions --- a critical path through the causal DAG --- and each step
must build on the last.

This imposes a physical requirement on the environment: it must maintain a
universally agreed-upon record of which interactions have occurred and in
which order. We call this record the \emph{causal ledger}. Environments in
which the causal ledger can be stably maintained permit biological complexity;
environments in which it fragments do not. This is not a probabilistic
statement but a structural one: a molecular assembly process requiring $k$
sequential irreversible steps cannot be completed in an environment where
different regions cannot agree on a causal ordering of events.

Two conditions are necessary for causal ledger preservation:

\begin{description}
  \item[Condition F1 (Quantum Darwinism redundancy).]
  The actualized state of each local interaction must be redundantly encoded
  in at least $R_{\min}$ independent environmental fragments. Following
  Zurek's Quantum Darwinism framework \citep{Zurek2009}, the redundancy
  $R := |\mathcal{E}|/I_\delta$ must satisfy $R \geq R_{\min}$, where
  $|\mathcal{E}|$ is the number of independent environmental fragments and
  $I_\delta$ is the partial information threshold. In practice
  $R_{\min} \sim \mathcal{O}(10)$ for macroscopic classical facts to emerge.
  Without this redundancy, actualized states cannot become objective classical
  facts, and sequential assembly steps cannot build reliably on one another.

  \item[Condition F2 (Non-relativistic causal regime).]
  The local proper acceleration $a$ must satisfy
  \begin{equation}
    T_U = \frac{\hbar a}{2\pi c k_B} \;\ll\; T_{\mathrm{env}},
    \label{eq:unruh}
  \end{equation}
  where $T_U$ is the Unruh temperature \citep{Unruh1976} and $T_{\mathrm{env}}$
  is the environmental temperature. In highly accelerated or strongly curved
  spacetime regions, Rindler observers perceive a thermal bath with temperature
  $T_U$. When $T_U \gtrsim T_{\mathrm{env}}$, the frame-dependence of event
  ordering disrupts the causal ledger: different observers cannot agree on the
  chronological sequence of interactions, and molecular assembly requiring a
  fixed causal order structurally fails.

  \paragraph{F2 as a bound on molecular machinery kinematics.}
  Condition F2 is not merely a cosmological constraint.  It bounds the
  \emph{kinematics of molecular machinery} during individual assembly events.
  In extreme prebiotic environments --- violent hydrothermal vent gradients,
  impact-shocked prebiotic soups, lightning-driven chemistry in reducing
  atmospheres --- molecular components undergo transient accelerations
  large on microscopic timescales.  F2 requires that even during these
  transient events, the Unruh temperature associated with the instantaneous
  acceleration of a reacting molecular unit must remain well below the local
  thermal background.  If it does not, the causal ordering of bond-forming steps
  is disrupted at the level of the individual reaction event, preventing the
  sequential accumulation of causal depth required by Condition~F1.

  For a molecule undergoing acceleration $a$, the numerical threshold is
  $T_U \approx 4\times10^{-24}\,a\,[\mathrm{m\,s}^{-2}]\,\mathrm{K}$.
  Even at $a = 10^{15}\,\mathrm{m\,s}^{-2}$ (extreme shock acceleration),
  $T_U \approx 4\times10^{-9}\,\mathrm{K}$, negligible against any chemically
  relevant $T_{\mathrm{env}}$.  F2 is therefore trivially satisfied in all
  terrestrial and most astrophysical chemistry, including the extreme environments
  most often proposed as origin-of-life sites.  Its scientific value is as a
  \emph{universal boundary condition}: it precisely identifies the regime --- near
  neutron star surfaces, within intense gravitational wave backgrounds, or during
  Planck-density events in the early universe --- where the causal preconditions
  for biological assembly are physically absent regardless of chemistry, and it
  does so without any appeal to habitability proxies.
\end{description}

\subsection{Substrate-Independence}

The central feature of the Causal Firewall for astrobiology is what it
\emph{does not} require. Neither F1 nor F2 makes any reference to:
\begin{itemize}[noitemsep]
  \item Molecular class (carbon, silicon, or any other element);
  \item Solvent chemistry (water, ammonia, liquid methane, or any other medium);
  \item Temperature range (beyond the implicit requirement that $T_U \ll T_{\mathrm{env}}$,
    which constrains acceleration, not temperature per se);
  \item Biological alphabet (nucleic acids, proteins, or any specific polymer class).
\end{itemize}
Condition F1 requires that the environment support redundant information
encoding in \emph{some} physical medium --- it does not specify what that
medium is. Condition F2 requires that $T_U/T_{\mathrm{env}} \ll 1$ --- a
ratio of temperatures that depends only on the local gravitational environment
and the ambient radiation field, not on chemistry. This substrate-independence
is not an assumption of the framework; it follows from the fact that the
Causal Firewall is defined at the level of the causal graph, which is logically
prior to any specification of chemistry.

\begin{remark}
The liquid water requirement of the CHZ framework is not wrong as a
terrestrial proxy; it is incomplete as a universal criterion. Liquid water
satisfies Conditions F1 and F2 in Earth's environment, which is why it
correlates with life on Earth. But the physical reason life exists on
Earth is not that water is present; it is that the causal ledger can be
maintained. Any environment in which the ledger is maintained is, in
principle, biologically permissive under RA --- regardless of whether
water is present.
\end{remark}

% ── 3. The μ = 1 Percolation Threshold ───────────────────────────────────────
\section{The Causal Firewall as a Percolation Threshold}
\label{sec:percolation}

The Causal Firewall has a precise mathematical formulation as a phase
transition in the local causal graph. Define:
\begin{itemize}[noitemsep]
  \item $\lambda(x,t)$: the local actualization rate density (interactions
    per unit 4-volume crossing the irreversibility threshold);
  \item $\tau_d$: the decoherence timescale (time for Quantum Darwinism
    redundancy $R$ to exceed $R_{\min}$);
  \item $\ell$: the characteristic coherence length of the molecular system.
\end{itemize}

The mean-field parameter of the local causal graph is
$\mu = \lambda\,\tau_d\,\ell^3$, which counts the expected number of new
irreversible interactions within one coherence volume per decoherence time.
By the Erd\H{o}s-R\'{e}nyi theorem for directed random graphs
\citep{ErdosRenyi1960}:

\begin{theorem}[Causal Firewall Phase Transition]
\label{thm:percolation}
The local causal graph has:
\begin{itemize}[noitemsep]
  \item $\mu < 1$: all strongly connected components are $\mathcal{O}(\ln N)$
    --- the graph fragments into trees, no universal causal record exists,
    biological complexity is structurally impossible.
  \item $\mu = 1$: critical point; the largest strongly connected component
    scales as $N^{2/3}$ --- the Causal Firewall threshold.
  \item $\mu > 1$: a giant strongly connected component of size
    $\approx(1-e^{-\mu})N$ emerges --- a shared causal ledger exists and
    biological complexity is possible.
\end{itemize}
The Causal Firewall condition is $\mu = \lambda\,\tau_d\,\ell^3 \geq 1$.
\end{theorem}

The shared causal ledger required for sequential molecular assembly is
precisely the giant strongly connected component described in
Theorem~\ref{thm:percolation}: the region in which every pair of vertices
has a directed causal path in both directions, ensuring that all observers
reconstruct the same actualization history \citep{Sandeman2026RACI}.

\paragraph{The universal $\mu = 1$ transition.}
The same Erd\H{o}s-R\'{e}nyi transition governs four apparently unrelated
physical phenomena in the RA framework:
\begin{itemize}[noitemsep]
  \item \textbf{QCD confinement}: quarks and gluons cannot propagate freely
    as asymptotic states because their causal topology reaches the BDG
    filter boundary at $\mu < 1$ \citep{Sandeman2026RATM}.
  \item \textbf{Galactic rotation curves}: the same $\mu = 1$ transition is
    proposed in RAGC to govern the flat rotation curve onset as the transition
    from $\mu > 1$ (disk) to $\mu < 1$ (halo) in the actualization density
    field~\citep{Sandeman2026RAGC} (a conjectural result pending full derivation
    of the covariant RA field equations in the galactic regime).
  \item \textbf{Quantum fault-tolerance}: the threshold for topological
    quantum error correction is the Kinematic Coherence Bound
    $N_{\max} = \eta p_{\mathrm{th}}$, determined by the same $\mu = 1$
    condition \citep{Sandeman2026RAQI}.
  \item \textbf{The Causal Firewall}: biological complexity requires
    $\mu \geq 1$ in the local causal graph, as shown above.
\end{itemize}
This four-way unification --- from particle physics to cosmology to quantum
computing to astrobiology --- is not a coincidence but a structural
consequence of the causal graph's Erd\H{o}s-R\'{e}nyi combinatorics.
Wherever the same mathematical transition appears, the same physical
phenomenon occurs: the emergence or suppression of coherent causal structure.

% ── 4. Assembly Depth and the Origin of Life ─────────────────────────────────
\section{Causal Assembly Depth and the Origin of Life}
\label{sec:ara}

\subsection{Coarse-Grained Assembly Depth}

The Assembly Theory framework of \citet{Murray2023} quantifies molecular
complexity by the assembly index $A(M)$ --- the minimum number of JOIN
operations required to construct molecule $M$ from elementary constituents.
In the RA framework, the \emph{Actualism Assembly Index} $\ARA(M) :=
\mathrm{depth}(G_M)$ gives the causal depth of the molecular subgraph,
with $\ARA(M) \geq A(M)$ (equality at the minimal assembly DAG)
\citep{Sandeman2026RAHC}.

For metabolic networks, a tractable network-level generalisation is the
\emph{coarse-grained assembly depth}:
\begin{equation}
  \tARA(\mathcal{N}) \;:=\; \text{number of irreversible reactions on
  the critical path of network } \mathcal{N},
  \label{eq:tara}
\end{equation}
where the critical path is the longest directed chain of reactions satisfying
the irreversibility condition $\Delta G < 0$ under physiological conditions
\citep{Sandeman2026RAHC}. This is computationally tractable from standard
metabolic databases (KEGG, BiGG) using longest-path algorithms on the
reaction DAG.

For glycolysis, the canonical 10-step pathway from glucose to pyruvate,
$\tARA = 3$: only three reactions are strongly irreversible under physiological
conditions (hexokinase: $\Delta G \approx -16.7$~kJ/mol; phosphofructokinase-1:
$\Delta G \approx -14.2$~kJ/mol; pyruvate kinase: $\Delta G \approx
-31.4$~kJ/mol). The critical path of length 3 is precisely the bottleneck
of the pathway, and allosteric regulation concentrates at exactly these
three nodes --- independent biological confirmation of the framework.

\subsection{The Origin-of-Life Sandwich Bound}

The Causal Firewall imposes a necessary condition on the first proto-metabolic
networks. A network whose critical path exceeds the available cycle time cannot
maintain $\mu \geq 1$ and therefore cannot sustain the causal coherence required
for replication. Combined with the Assembly Theory lower bound, this gives:

\begin{proposition}[Sandwich bound on the first replicators]
\label{prop:sandwich}
The proto-metabolic network $\mathcal{N}_0$ of the first self-replicating
systems must satisfy:
\begin{equation}
  A(M^*_0) \;\leq\; \tARA(\mathcal{N}_0)
  \;\leq\; \frac{\tau_{\mathrm{cycle},0}}{\tau_{\mathrm{irrev}}},
  \label{eq:sandwich}
\end{equation}
where $M^*_0$ is the most assembly-complex essential product, $\tau_{\mathrm{cycle},0}$
is the proto-replication cycle time, and $\tau_{\mathrm{irrev}}$ is the mean
time per irreversible reaction step.
\end{proposition}

The left inequality follows from the Assembly Index Correspondence
\citep{Sandeman2026RAHC}: a network cannot produce a molecule more complex
than its Assembly Index without a causal path of at least that depth.
The right inequality follows from the Causal Firewall compatibility condition
(Proposition~2.2 of \citet{Sandeman2026RACI}):
a cell cycle cannot complete faster than its critical path allows.

This double bound is a \emph{quantitative} constraint on origin-of-life models,
not merely a qualitative one. Given measured assembly indices for candidate
first replicators (RNA monomers, peptides, autocatalytic lipids) and estimated
early-Earth reaction timescales, the sandwich~\eqref{eq:sandwich} is
evaluable without new experimental data. A proposed first replicator is
physically inadmissible under RA if $A(M^*_0)$ exceeds a realistic
$\tau_{\mathrm{cycle},0}/\tau_{\mathrm{irrev}}$ for early-Earth conditions.

% ── 5. Universal Biosignature Criteria ───────────────────────────────────────
\section{Universal Biosignature Criteria and SETI Implications}
\label{sec:predictions}

The Causal Firewall conditions translate directly into three astrophysical
observables, each measurable for any candidate environment independently
of its chemistry.

\subsection{Criterion 1: The Unruh Temperature Ratio}

The Unruh temperature $T_U = \hbar a / (2\pi c k_B)$ is calculable from
the local proper acceleration $a$, which is directly related to the
surface gravity of a body or the curvature of the local spacetime.
For a planetary surface ($g = 9.8$~m/s$^2$):
$T_U \approx 4 \times 10^{-20}$~K --- negligible against any environmental
temperature, consistent with life being possible at planetary surfaces.
The criterion becomes discriminating at proper accelerations of order
$a_c \sim 10^{20}$~m/s$^2$, at which $T_U \sim 1$~K and begins to compete
with environmental temperatures in cold astrophysical environments. Such
accelerations are found in relativistic jets ($a \sim 10^{22}$--$10^{25}$~m/s$^2$
for plasma in the jet frame), in the innermost stable circular orbits of stellar
black holes, and in the tidal fields of neutron star binary mergers --- not at
neutron star surfaces, where $T_U \sim 4 \times 10^{-9}$~K remains safely
below the $\sim 10^6$~K surface temperature. The Causal Firewall predicts
that wherever $T_U / T_{\mathrm{env}} \gtrsim 0.1$, the causal ledger is
sufficiently disrupted that complex molecular assembly is structurally
suppressed.

\subsection{Criterion 2: Decoherence Redundancy}

The Quantum Darwinism condition (F1) requires a matter-dense environment
with sufficient molecular scattering cross-section to produce $R \geq
R_{\min} \sim 10$ independent environmental records of each interaction.
This condition is correlated with --- but not identical to --- matter
density. High-energy, low-density environments (the intergalactic medium,
inner radiation belts) may have high interaction rates but insufficient
environmental complexity for redundant encoding. The relevant observable
is the product of matter density and scattering cross-section per unit
volume, which is in principle estimable from spectroscopic data for
any candidate environment.

\subsection{Criterion 3: Causal Coherence Density}

The condition $\mu = \lambda\,\tau_d\,\ell^3 \approx 1$ requires sufficient
actualization flux density to maintain the giant causal component without
overwhelming it. This is a Goldilocks condition on the thermodynamic flux
through the local environment: enough far-from-equilibrium chemistry to
drive the causal graph past criticality, but not so much that causal
severances fragment the ledger faster than it can be written. Environments
satisfying this condition are precisely those with moderate matter density,
persistent thermodynamic gradients, and low proper acceleration: the
description of planetary surfaces around stable main-sequence stars.

\subsection{Falsifiable Predictions}

Two predictions follow from the framework:

\begin{enumerate}
  \item \textbf{Life cannot exist in environments where $T_U \gtrsim
  T_{\mathrm{env}}$.} This excludes environments near neutron stars, in
  relativistic jets, and in highly curved spacetime regions, regardless
  of energy availability or chemical composition. The prediction is
  falsifiable: discovery of life in an environment demonstrably satisfying
  $T_U \gtrsim T_{\mathrm{env}}$ would require fundamental revision of
  the framework.

  \item \textbf{All discovered extraterrestrial life satisfies Conditions
  F1 and F2.} This is not a tautology but a non-trivial constraint: it
  predicts that life will not be found in environments that satisfy water
  or energy criteria but violate the causal coherence conditions. The
  Europa subsurface ocean, the Enceladus plumes, and the Titan lakes are
  ordered differently by Causal Firewall criteria than by CHZ criteria.
  RA predicts that Causal Firewall-satisfying environments are the correct
  predictor of habitability, not CHZ membership.
\end{enumerate}

% ── 6. Comparison with Existing Frameworks ───────────────────────────────────
\section{Relation to Existing Habitability Frameworks}
\label{sec:comparison}

\paragraph{The circumstellar habitable zone.}
The CHZ \citep{Kasting1993, Kopparapu2013} is a geometrically defined
region around a star where a planet could maintain surface liquid water.
It is an excellent practical tool for prioritising telescope observations,
but it is a proxy for habitability rather than a physical criterion for
it. The Causal Firewall explains \emph{why} the CHZ works as a proxy on
Earth: planetary surfaces in the CHZ tend to satisfy Conditions F1 and F2.
But the proxy breaks down at the edges: high-energy bodies within the CHZ
may violate F2 (excessive curvature or relativistic motion of the host
system); subsurface environments outside the CHZ may satisfy both conditions
(Europa, Enceladus). The Causal Firewall criterion is more discriminating
in some directions and more permissive in others.

\paragraph{Alternative biochemistry.}
Proposals for non-aqueous life \citep{Benner2004} expand the chemical
palette but do not identify the physical ground of habitability. The RA
framework provides that ground: non-carbon, non-water life is not prohibited,
but it is subject to the same Causal Firewall conditions as carbon-water
life. The question for any proposed alternative biochemistry is not ``can
this chemistry support complex molecules?'' but ``can this chemistry
maintain a causal ledger?'' The answer depends on the physical environment,
not on the molecular class.

\paragraph{Dissipative adaptation.}
England's dissipative adaptation framework \citep{England2013} predicts
that matter driven far from equilibrium by thermodynamic flux will tend
to evolve configurations that dissipate energy more effectively. This
is a kinetic argument for the emergence of self-replicating chemistry.
In RA, this argument is strengthened: the replication rate is
exponentially suppressed by the causal depth $\ARA(M)$ of the critical
molecules, so selection for replication speed is simultaneously selection
for minimal causal assembly depth. The Causal Firewall provides the
structural condition under which England's kinetic process can operate;
the sandwich bound provides a quantitative constraint on its first
products.

\paragraph{Assembly Theory.}
The Assembly Theory framework \citep{Murray2023} provides an empirical
tool for measuring molecular complexity via the assembly index $A(M)$,
and has been proposed as a life-detection criterion: high-$A$ molecules
are unlikely to be produced abiotically. The RA framework provides the
physical grounding for this connection: $A(M) \leq \ARA(M)$, with equality
at the minimal causal path, and the assembly index is a lower bound on
the causal depth of production. The Causal Firewall explains why high-$A$
molecules are biosignatures: they require a causal ledger-preserving
environment to be produced, and such environments are precisely the
ones where life is possible.

% ── 7. Conclusion ─────────────────────────────────────────────────────────────
\section{Conclusion}
\label{sec:conclusion}

We have derived a substrate-independent physical criterion for the
distribution of biological complexity in the universe. The Causal Firewall
--- two conditions on the local causal graph structure, neither of which
makes reference to chemistry --- provides a more fundamental habitability
criterion than the circumstellar habitable zone, and one that is directly
measurable as astrophysical observables.

The key results are:
\begin{enumerate}[noitemsep]
  \item The Causal Firewall corresponds to the $\mu = 1$ Erd\H{o}s-R\'{e}nyi
    percolation transition --- the same mathematical threshold governing QCD
    confinement, galactic rotation curves, and quantum fault-tolerance.
  \item Non-carbon-based life is not prohibited but is subject to the same
    causal coherence requirements as carbon-water life.
  \item The sandwich bound~\eqref{eq:sandwich} provides quantitative
    constraints on the first self-replicating systems evaluable with existing
    Assembly Theory tools.
  \item SETI search programmes are most effectively weighted toward environments
    satisfying Conditions F1 and F2 --- not merely toward energy-rich or
    chemically favourable environments.
\end{enumerate}

The framework makes two falsifiable predictions: life does not exist where
$T_U \gtrsim T_{\mathrm{env}}$, and all discovered extraterrestrial life
satisfies both Causal Firewall conditions. These predictions will become
testable as the search for extraterrestrial life matures. The programme of
measurement they require --- the Unruh temperature ratio, the decoherence
redundancy, and the causal coherence density of candidate environments ---
is physically motivated and substrate-independent.

\section*{Acknowledgements}

The author is an independent researcher with no institutional affiliation.
This research received no external funding. The author declares no
conflicts of interest.

\smallskip\noindent
\textbf{AI use disclosure.}
Claude (Anthropic, model \texttt{claude-sonnet-4-6}) was used in the
preparation of this manuscript, including drafting and revision of text,
\LaTeX{} compilation, and mathematical verification scaffolding.
The author is solely responsible for all scientific content,
mathematical claims, and conclusions.
AI tools were used in accordance with Cambridge University Press's
generative AI policy.

\smallskip\noindent
\textbf{Data availability.}
No new data were created or analysed. All supporting derivations are
available in the companion papers cited below and archived at
\url{https://github.com/jsandeman/Relational-Actualism}.

\begin{thebibliography}{99}

\bibitem[Benner et~al.(2004)]{Benner2004}
S.~A. Benner, A.~Ricardo, and M.~A. Carrigan.
\newblock Is there a common chemical model for life in the universe?
\newblock \textit{Curr.\ Opin.\ Chem.\ Biol.}, 8:672--689, 2004.

\bibitem[England(2013)]{England2013}
J.~L. England.
\newblock Statistical physics of self-replication.
\newblock \textit{J.\ Chem.\ Phys.}, 139:121923, 2013.

\bibitem[Erd\H{o}s \& R\'{e}nyi(1960)]{ErdosRenyi1960}
P.~Erd\H{o}s and A.~R\'{e}nyi.
\newblock On the evolution of random graphs.
\newblock \textit{Publ.\ Math.\ Inst.\ Hung.\ Acad.\ Sci.}, 5:17--61, 1960.

\bibitem[Kasting et~al.(1993)]{Kasting1993}
J.~F. Kasting, D.~P. Whitmire, and R.~T. Reynolds.
\newblock Habitable zones around main sequence stars.
\newblock \textit{Icarus}, 101:108--128, 1993.

\bibitem[Kopparapu et~al.(2013)]{Kopparapu2013}
R.~K. Kopparapu et al.
\newblock Habitable zones around main-sequence stars: new estimates.
\newblock \textit{Astrophys.\ J.}, 765:131, 2013.

\bibitem[Murray et~al.(2023)]{Murray2023}
J.~Murray et al.
\newblock Assembly theory explains and quantifies selection and evolution.
\newblock \textit{Nature}, 622:321--328, 2023.

\bibitem[Sandeman(2026a)]{Sandeman2026RAQM}
J.~F. Sandeman.
\newblock Relational Actualism and Quantum Mechanics (RAQM).
\newblock \textit{Foundations of Physics} (under review), 2026.
\newblock \href{https://doi.org/10.5281/zenodo.19174380}{doi:10.5281/zenodo.19174380}.

\bibitem[Sandeman(2026b)]{Sandeman2026RAGC}
J.~F. Sandeman.
\newblock Relational Actualism: Gravity and Cosmology (RAGC).
\newblock \textit{Preprint}, 2026.
\newblock \href{https://doi.org/10.5281/zenodo.19197999}{doi:10.5281/zenodo.19197999}.

\bibitem[Sandeman(2026c)]{Sandeman2026RAHC}
J.~F. Sandeman.
\newblock Relational Actualism: Algorithmic Causal Dynamics (RAHC).
\newblock \textit{Preprint}, 2026.
\newblock \href{https://doi.org/10.5281/zenodo.19198171}{doi:10.5281/zenodo.19198171}.

\bibitem[Sandeman(2026d)]{Sandeman2026RACI}
J.~F. Sandeman.
\newblock Relational Actualism: Complexity and Intelligence (RACI).
\newblock \textit{Preprint}, 2026.
\newblock \href{https://doi.org/10.5281/zenodo.19198224}{doi:10.5281/zenodo.19198224}.

\bibitem[Sandeman(2026e)]{Sandeman2026RATM}
J.~F. Sandeman.
\newblock Relational Actualism: The Topology of Matter (RATM).
\newblock \textit{Preprint}, 2026.
\newblock \href{https://doi.org/10.5281/zenodo.19265651}{doi:10.5281/zenodo.19265651}.

\bibitem[Sandeman(2026f)]{Sandeman2026RAQI}
J.~F. Sandeman.
\newblock Relational Actualism: Implications for Quantum Information (RAQI).
\newblock \textit{Preprint}, 2026.
\newblock \href{https://doi.org/10.5281/zenodo.19198285}{doi:10.5281/zenodo.19198285}.

\bibitem[Unruh(1976)]{Unruh1976}
W.~G. Unruh.
\newblock Notes on black-hole evaporation.
\newblock \textit{Phys.\ Rev.\ D}, 14:870--892, 1976.

\bibitem[Zurek(2009)]{Zurek2009}
W.~H. Zurek.
\newblock Quantum Darwinism.
\newblock \textit{Nature Phys.}, 5:181--188, 2009.

\end{thebibliography}

\end{document}
